
\section{Introduction to Capstone}
\label{capstone}

\textbf{Quick Reference Guide}

\begin{center} 
\begin{tabular}{|p{1.8in}|p{2.5in}|c|} 
\hline \centering{\textbf{What You Want To Do}} & \centering{\textbf{How You Do It}} & \textbf{Button} \\ 

\hline Start Recording Data 
& Click the `Record' button
& \raisebox{-5mm}{\includegraphics{../../131/StudentGuideModule1/appendices/capstone/record_button.eps}} \\ 

\hline Stop Recording Data 
& Click the `Stop' button
& \raisebox{-5mm}{\includegraphics{../../131/StudentGuideModule1/appendices/capstone/stop_button.eps}} \\ 

\hline Re-scale the data so it fills the Graph display window 
& Click the `Scale to Fit' button
& \raisebox{-5mm}{\includegraphics{../../131/StudentGuideModule1/appendices/capstone/scale_to_fit.eps}} \\ 

\hline Pinpoint the x- and y-coordinate values on the Graph display
& Click the `Delta Tool' button, move square to point of interest
& \raisebox{-5mm}{\includegraphics{../../131/StudentGuideModule1/appendices/capstone/delta_tool.eps}} \\ 

\hline 
\end{tabular} 
\end{center}
\index{color page}

\bigskip
\textbf{Selecting a Section of Data}

\begin{enumerate}
\item To select a data section, click the \textbf{Highlight} button.  Drag the pink rectangle to the data you wish to highlight.  Adjust the width and height using horizontal and vertical arrows at edges of highlighted region.
\item To unselect the data, with cursor in highlighted region, right click and select ``Delete highlighter''.
\end{enumerate}
\textbf{Finding Average and Standard Deviation of a Section of Data}

\begin{enumerate}
\item Select the section of data of interest, as described above.
\item Click on the \textbf{Statistics} button on the Graph Toolbar. Use the menu on this button to select maximum, minimum, mean, standard deviation, or any combination.
\item To remove the results, click on the \textbf{Statistics} button again.
\end{enumerate}


\pagebreak[2]
\textbf{Finding the Area Under a Curve}

\begin{enumerate}
\item Select the section of data that you want to integrate under, using the procedure described above.
\item Click the \textbf{Display Area} button on the Graph Toolbar.
The results of the integration will be displayed on the graph.
\item To undo the integration, click on the \textbf{Display Area} button again.
\end{enumerate}

\textbf{Calibrating Force Sensors}

\begin{enumerate}
\item Connect force sensor to Pasco interface (in port ``A'').

\item Open whatever software application you will be using in the current experiment.

\item Click ``Calibration'' icon on the tools panel on the left side of the screen.

\item Verify that the type of measurement is set to ``Force'', then click ``Next''.

\item Choose ``Two standards'' as calibration type, then click ``Next''.

\vspace{0.1in}
\hspace{-0.2in}FOR HORIZONTAL USE:

\item Set force sensor on track, and with no force applied to the sensor, press TARE button on side of the sensor. This sets the sensor to zero.

\item With no mass attached, set standard value to 0 for the first point, click ``set current value to standard value'', then click ``Next''.

\item While holding sensor stationary (with your hand), hang 200g from sensor over a pulley at end of track, as in Figure 2 of the Force and Motion I experiment. Set standard value to 1.96 newtons, click ``set current value to standard value'', then click ``Next''.

\item After reviewing calibration, click ``Finish''.

\item Click ``Calibration'' button again to close. Force sensor is now calibrated for the rest of your experiment.

\item While still holding the force sensor still, press ``Record''. Graph should now show a reading of 1.96 N. Press ``Stop''.

\item Try hanging a different mass from the force sensor (over the pulley); press ``Record'' and check that it is reading correctly.

\vspace{0.1in}
\hspace{-0.2in}FOR VERTICAL USE:

\setcounter{enumi}{4}

\item Holding sensor vertically (with hook down and no force applied to it), press TARE button on side of the sensor. This sets the sensor to zero.

\item With no mass attached, set standard value to 0 for the first point, click ``set current value to standard value'', then click ``Next''.

\item Hang 200 g from sensor, set standard value to 1.96 newtons, click ``set current value to standard value'', then click ``Next''.

\item After reviewing calibration, click ``Finish''.

\item Click ``Calibration button again to close. Force sensor is now calibrated for the rest of your experiment.

\item While still holding the force sensor still, press ``Record''. Graph should now show a reading of 1.96 N. Press ``Stop''.

\item Try hanging a different mass from the force sensor; press ``Record'' and check that it is reading correctly.
\end{enumerate}
